\documentclass{article}

\usepackage[utf8]{inputenc}  % Allows you to type UTF-8 characters directly
\usepackage[T1]{fontenc}    % Ensures the PDF output uses 8-bit encoding for better font rendering

\title{wspolrzedne-grafu}
\author{Ksawery Olaf }
\date{March 2026}

\begin{document}

\maketitle

\section{Dokumentacja funkcjonalna opisująca działanie projektu}

\subsection{Cel projektu}
Celem niniejszej aplikacji jest automatyczne wyznaczanie optymalnych współrzędnych wierzchołków dla grafów planarnych. Program przekształca surową listę krawędzi w czytelną reprezentację wizualną, wykorzystując zaawansowane algorytmy rozmieszczania wierzchołków (tzw. \textit{graph layout algorithms}), co pozwala na uniknięcie nakładania się krawędzi i zwiększenie czytelności struktury danych.

\subsection{Opis funkcji programu}
Program oferuje następujące funkcjonalności:
\begin{itemize}
    \item Wczytywanie struktury grafu z plików tekstowych (lista krawędzi).
    \item Obliczanie współrzędnych $x, y$ dla każdego węzła przy użyciu co najmniej dwóch algorytmów (np. Fruchterman–Reingold, Twierdzenie Tutte’a).
    \item Eksport wyników do formatu tekstowego lub binarnego, zgodnie z wyborem użytkownika.
    \item Wersja C: Szybkie przetwarzanie wsadowe za pomocą interfejsu linii poleceń (CLI).
    \item Wersja Java: Interaktywna wizualizacja grafu z możliwością ręcznej modyfikacji widoku.
\end{itemize}

\subsection{Interfejs użytkownika i obsługa programu}

\subsubsection{Aplikacja w języku C (CLI)}
Program sterowany jest za pomocą flag i argumentów wywołania.
\begin{description}
    \item[\texttt{-i <plik>}] Określa plik wejściowy z listą krawędzi.
    \item[\texttt{-o <plik>}] Określa plik wyjściowy dla współrzędnych.
    \item[\texttt{-a <algorytm>}] Wybór algorytmu obliczeniowego (np. \texttt{FR} dla Fruchterman-Reingold, \texttt{TUT} dla Tutte).
    \item[\texttt{-b}] Zapis wyniku w formacie binarnym (domyślnie tekstowy).
\end{description}

\subsubsection{Aplikacja w języku Java (GUI)}
Interfejs graficzny umożliwia:
\begin{itemize}
    \item \textbf{Menu Plik:} Wczytywanie danych (tekstowych/binarnych) oraz zapis wyników.
    \item \textbf{Panel Wizualizacji:} Samodzielnie zaimplementowany komponent rysujący graf.
    \item \textbf{Panel Narzędziowy:} Wybór algorytmu, zmiana pola widzenia (zoom), oraz przełączanie etykiet/wag.
\end{itemize}

\subsection{Format danych wejściowych i wyjściowych}
\paragraph{Dane wejściowe (TXT):} Plik powinien zawierać wiersze w formacie: \\
\texttt{<nazwa\_krawędzi> <wierzchołek\_A> <wierzchołek\_B> <waga>}.

\paragraph{Dane wyjściowe (TXT):} Wynikiem jest lista współrzędnych: \\
\texttt{<id\_wierzchołka> <x> <y>}.

\subsection{Sytuacje wyjątkowe i komunikaty błędów}
Program przewiduje i obsługuje następujące błędy:
\begin{table}[h]
\centering
\begin{tabular}{|l|l|}
\hline
\textbf{Sytuacja} & \textbf{Komunikat / Zachowanie} \\ \hline
Brak pliku wejściowego & "Error: Input file not found." (Kod powrotu 1) \\ \hline
Nieprawidłowy format danych & "Error: Invalid graph format at line [nr]." \\ \hline
Graf nie jest planarny & "Warning: Provided graph may not be planar." \\ \hline
Błąd zapisu & "Error: Permission denied / Disk full." \\ \hline
\end{tabular}
\caption{Obsługa błędów w aplikacji.}
\end{table}
\section{Dokumentacja implementacyjna}
\section{Końcowa dokumentacja projektu}

\end{document}
